\documentclass[11pt]{article}
\usepackage[margin=1in]{geometry}
\usepackage{amsmath,amssymb,amsfonts,mathtools,bm}
\usepackage[T1]{fontenc}
\usepackage[utf8]{inputenc}
\title{Inline and Display Mathematics}
\author{Source-backed grouped formula sample}
\date{}
\begin{document}
\maketitle
\section{Expressions}
\paragraph{Expression 1.} The following expression is taken from a source-backed formula corpus.

\begin{align*}\gamma_{g,9}:|D_{\rm inv}\rangle \rightarrow {1\over\sqrt{D}} \sum_\alpha \sum_{\beta,\gamma} c_{\alpha\beta} {\overline c}_{\gamma\alpha} |d_\beta\otimes{\overline d}_\gamma \rangle={1\over \sqrt{D}} \sum_\alpha |d_\alpha\otimes{\overline d}_\alpha\rangle=|D_{\rm inv}\rangle~,\end{align*}

\paragraph{Expression 2.} The following expression is taken from a source-backed formula corpus.

\begin{align*}\lim_{r\to\infty}\int_{{\cal I}^-} dt\, d\phi\, dt'\, d\phi'\,\left[\frac{(rr')^2}{l^2}\left(\Phi(r,t,\phi){\buildrel \leftrightarrow \over \partial_{r_{\ast}}}G(r,t,\phi; r',t',\phi'){\buildrel \leftrightarrow \over \partial_{r'_{\ast}}}\Phi(r',t',\phi')\right)\right]_{r=r'}, \end{align*}

\paragraph{Expression 3.} The following expression is taken from a source-backed formula corpus.

\begin{align*}\langle E_{eff}(x_0,x_1) \rangle \;=\; \frac{1}{6} \left[\frac{\partial^2}{\partial x_0^2}\frac{\partial^2}{\partial y_1^2} -\frac{\partial^2}{\partial x_0 \partial x_1} \frac{\partial^2}{\partial y_0 \partial y_1}\right] \Delta_\varphi(x_0,x_1;y_0, y_1)|_{y_\mu \to x_\mu} \;,\end{align*}

\paragraph{Expression 4.} The following expression is taken from a source-backed formula corpus.

\begin{align*}\rule{.9cm}{0cm} \frac{\delta}{\delta\phi^*_A}+ \frac{i}{\hbar}\frac{\delta\Gamma}{\delta\phi^*_A} -(-1)^{\varepsilon_i(\varepsilon_A+1)}(\Gamma^{''-1})^{ij} \bigg(\frac{\delta_l}{\delta\varphi^j}\frac{\delta\Gamma}{\delta\phi^*_A} \bigg)\frac{\delta_l}{\delta\varphi^i}\bigg),\end{align*}

\paragraph{Expression 5.} The following expression is taken from a source-backed formula corpus.

\begin{align*}\hat h_{\rm red} = \sum_{\vec{b}}\left\{-\frac{g^2 LM}{4\sqrt2}\frac{1}{J}\frac{\delta}{\delta\varphi_{\vec b}} J\frac{\delta}{\delta\varphi _{\vec b}}+\left(\frac{4\sqrt2}{g^2LM}\right) \sum_{\vec{\varepsilon}}(\varphi_{\vec b}-\varphi_{\vec b+\vec\varepsilon})^2\right\} \, .\end{align*}
\end{document}
